\documentclass[12pt]{article}
\usepackage[utf8]{inputenc}
\usepackage[T1]{fontenc}
\usepackage{hyperref}
\title{Eastern Franconian Language Resources for K-12 Curriculum Development}
\author{Compiled by Claude Sonnet 5 (claude-sonnet-5) for Andr\'as Kornai}
\date{}
\begin{document}
\maketitle
\section{Introduction}
Eastern Franconian (ISO 639-3 vmf) is a group of Upper German dialects of northern Bavaria and neighbouring areas of Germany. This search found no school curricula or teaching materials in the dialect. What turned up is mostly lexicographic and dialectological. Search coverage was limited to one query.
\section{Primary Educational Resources}
No primers, textbooks or school programmes were found in the search results.
\section{Linguistic and Reference Materials}
\begin{itemize}
\item Wikipedia overview of East Franconian German: \url{https://en.wikipedia.org/wiki/East_Franconian_German}
\item Bavarian Academy of Sciences article on the Fr\"ankisches W\"orterbuch project (about 6 million citations): \url{https://badw.de/fileadmin/pub/akademieAktuell/2007/23/03_Klepsch.pdf}
\item Lexilogos list of older Eastern Franconian dictionaries (Sartorius 1862 Würzburg, Heilig 1894 Tauber region, Spie{\ss} 1881 Henneberg): \url{https://www.lexilogos.com/english/franconian_east_dictionary.htm}
\item Digitized Würzburg dialect dictionary at the Deutsche Digitale Bibliothek: \url{https://www.deutsche-digitale-bibliothek.de/item/3QIFN2W6LSR22MOUQ4PTKBESNDA6I4UB}
\item Review of a Rothenburg dialect dictionary (over 3,000 entries, 1973): \url{https://journals.wlb-stuttgart.de/index.php/wfr/article/download/11450/11320/21680}
\end{itemize}
\section{Religious and Community Resources}
None found in the search results.
\section{Digital and Audio Resources}
A collection of poems in Franconian dialect is listed at the Deutsche Digitale Bibliothek (content not inspected): \url{https://www.deutsche-digitale-bibliothek.de/item/2XQQO3OHR66RSAX6JWQLMFWSZ3YUWY3S}
\section{Curriculum Development Context}
The dialect is well documented lexicographically, but no classroom materials were found. A curriculum would have to be built from the dictionaries and dialect studies above. Further German-language searches (for example on Mundart in Franken schools) are recommended.
\end{document}